\documentclass{training}
\setlang{french}
% Filigrane PAR MODULE : \uxfond (PNG A4 dilue 30 %). Cadre noir vectoriel commun.
% Proportions original 744 px : marge 20/744, trait 1/744. Exterieur BLANC.
% Cadre via eso-pic +\put (coordonnees absolues) — pas de remember picture (decalage).
\usepackage{eso-pic}

\providecommand{\uxfond}{../../resources/fond/norme_fond_paliers_30.png}

\newlength{\uxmargin}
\newlength{\uxthick}
\newlength{\uxinnerw}
\newlength{\uxinnerh}
\setlength{\uxmargin}{0.02688\paperwidth}% 20/744
\setlength{\uxthick}{0.001344\paperwidth}% 1/744
\setlength{\uxinnerw}{\dimexpr\paperwidth-2\uxmargin-2\uxthick\relax}
\setlength{\uxinnerh}{\dimexpr\paperheight-2\uxmargin-2\uxthick\relax}

% Contenu strictement a l'interieur du trait (+ jour confortable).
\geometry{%
  includeheadfoot,
  left=22mm,
  right=22mm,
  top=20mm,
  bottom=18mm,
  headheight=30pt,
  headsep=5mm,
  footskip=12mm%
}

\AddToShipoutPictureBG{%
  \AtPageLowerLeft{%
    % Fond blanc page entiere
    \put(0,0){\color{white}\rule{\paperwidth}{\paperheight}}%
    % Filigrane module a l'interieur du cadre
    \put(\LenToUnit{\dimexpr\uxmargin+\uxthick\relax},\LenToUnit{\dimexpr\uxmargin+\uxthick\relax}){%
      \includegraphics[width=\uxinnerw,height=\uxinnerh]{\uxfond}}%
    % Cadre noir : bas, haut, gauche, droite
    \put(\LenToUnit{\uxmargin},\LenToUnit{\uxmargin}){%
      \color{black}\rule{\dimexpr\paperwidth-2\uxmargin\relax}{\uxthick}}%
    \put(\LenToUnit{\uxmargin},\LenToUnit{\dimexpr\paperheight-\uxmargin-\uxthick\relax}){%
      \color{black}\rule{\dimexpr\paperwidth-2\uxmargin\relax}{\uxthick}}%
    \put(\LenToUnit{\uxmargin},\LenToUnit{\uxmargin}){%
      \color{black}\rule{\uxthick}{\dimexpr\paperheight-2\uxmargin\relax}}%
    \put(\LenToUnit{\dimexpr\paperwidth-\uxmargin-\uxthick\relax},\LenToUnit{\uxmargin}){%
      \color{black}\rule{\uxthick}{\dimexpr\paperheight-2\uxmargin\relax}}%
  }%
}

\makeatletter
\fancyhf{}
\fancyheadoffset{0pt}
\renewcommand{\headrulewidth}{0pt}
\lhead{\parbox[t]{0.56\textwidth}{\RaggedRight\footnotesize\sloppy\@title}}
\rhead{\parbox[t]{0.36\textwidth}{\RaggedLeft\footnotesize\sloppy\@subtitle}}
\cfoot{\thepage}
% Coupures : underscore et slash dans les noms de fichiers / consignes
\DeclareRobustCommand{\_}{\textunderscore\allowbreak}
\makeatother

% Evite les depassements de longues lignes (Github, Name_FirstName_...)
\sloppy
\emergencystretch=3em
\hbadness=10000

% Sommaire type 42 : chapitres seulement
\setcounter{tocdepth}{0}
\renewcommand{\cftchappresnum}{}
\renewcommand{\cftchapaftersnum}{}
\setlength{\cftchapnumwidth}{2.0em}
\setlength{\cftbeforechapskip}{0.35em}
\cftsetindents{chapter}{0em}{2.0em}


\begin{document}

\schoollogo{../../resources/jeunecrack.png}
\title{Heuristiques}
\subtitle{Piscine UX/UI --- UI}
\member{Piscine UX/UI}{https://github.com/01-edu/UX-UI}
\summary{Un audit, puis des recommandations qui nomment le critère.}
\version{1.0}
\prerequisite{Le module précédent de cette piscine, ou aucun si c'est le premier.}
\maketitle

\newpage
\tableofcontents
\newpage


\chapter{Instructions}

L'avant-propos n'est pas un exercice. Il ne sera pas noté. Les exercices commencent au chapitre suivant, du contrat le plus simple au plus dur. Un exercice plus tardif n'est pas pris en compte si un plus tôt ne tient pas.

On ne regarde que ce qui est dans le répertoire de rendu de l'exercice. N'ajoutez pas un fichier que l'exercice ne nomme pas. Le nom d'un livrable suit le motif imprimé dans l'exercice : \texttt{Name\_FirstName\_DeliverableName\_Date\_VersionNumber}.

Le geste ne vous est pas donné. Les exemples d'un exercice portent parfois une contrainte que la phrase ne répète pas. Lisez-les. Vos autres références sont le manuel de l'outil choisi, la personne assise à côté, et les liens imprimés dans l'exercice.

L'ordre est celui de la quête. Sauter une étape parce qu'un écran plus tard semble plus intéressant produit un écran que personne ne sait expliquer.

\subsection*{Des images}

La marque en couverture est le logo du projet. L'image de l'avant-propos est l'illustration de ce module. Le fond pâle derrière le texte est le filigrane de page. Aucun n'est un indice.

\subsection*{Des machines qui écrivent}

Vous pouvez vous servir d'un outil pour vérifier un mot, comme d'un dictionnaire. Vous ne pouvez pas remplacer un entretien, un test ou un fichier par un texte généré et le signer. Si vous ne pouvez pas montrer la personne, l'écran ou la ligne que vous avancez, l'exercice n'est pas fait.


\chapter{Avant-propos}

Steve Krug a appelé le livre \emph{Don't Make Me Think} et l'a passé à regarder des gens échouer en silence. Le titre est la culture. L'audit est le vôtre. Une liste copiée d'une affiche de dix slogans, sans écran attaché, n'est pas l'exercice.

\begin{center}
\includegraphics[width=0.78\textwidth,height=0.32\textheight,keepaspectratio]{../../resources/intros/UI-04-heuristics.jpg}\\[0.4em]
{\footnotesize Image : Shixart1985 / CC BY 2.0}
\end{center}

À regarder. Les vidéos sont un contexte. Elles ne sont pas le livrable.
\begin{itemize}
\item UI - Introduction \& Quest n°1, 01Edu \url{https://www.youtube.com/watch?v=UbE6WMeuvQg&list=PLHyAJ_GrRtf-_aBft2jQL4-QAunn2bLvu}
\end{itemize}

\chapter{Exercice 00 : UI V - Ex 1 - Audit}
\exnumber{00}
\extitle{UI V - Ex 1 - Audit}
\exfiles{ex00/archive zip, nomme Name\_FirstName\_DeliverableName\_Date\_VersionNumber}
\exauthorize{Ceux nommés dans l'exercice}
\makeheaderfiles

\textbf{Contexte :}

Bienvenue dans cette toute dernière quête !

L'objectif est d'endosser le rôle d'auditeur UX. On vous demandera de réaliser l'audit d'un site web ou d'une application, puis de rédiger une série de recommandations pour l'améliorer.

Soyons critiques, créatifs et impactants !

\textbf{Instructions :}

Réalisez un audit heuristique sur l'un de ces 3 sites web ou applications différents :

\begin{itemize}
\item https://www.ikea.com/ \url{https://www.ikea.com/}
\item https://www.doctorswithoutborders.org/ \url{https://www.doctorswithoutborders.org/}
\item Application WikiArt
\end{itemize}

\textbf{Livrables :}

\begin{itemize}
\item Un graphique complété
\item Un rapport comprenant les 5 étapes recommandées décrites ci-dessous
\item Déposez les documents sur Github, dans un zip nommé « ProjectTitle\_Name\_FirstName », et chaque livrable ainsi : « Name\_FirstName\_DeliverableName\_Date\_VersionNumber ». Exemple : « Doe\_John\_MidFiPrototype\_05242024\_V1 ».
\item Ajoutez un titre dans tous les documents écrits.
\item Soignez l'aspect global des livrables. Ils doivent être clairs, simples et lisibles. Vous pouvez regarder canva.com, sans surcharger le dessin de détails.
\end{itemize}

\textbf{Pistes :}

\begin{itemize}
\item Étapes recommandées
\begin{itemize}
\item Définir le périmètre.
\end{itemize}
\end{itemize}
1. Connaître les exigences business et le profil démographique des utilisateurs finaux.
2. Décidez quels outils de reporting et quelles heuristiques utiliser (dans ce cas : cette checklist d'évaluation heuristique \url{https://drive.google.com/file/d/10KbfbNZA1oVS1sXbjjXLPPmdZ6nqVkdc/view}).
3. Évaluer l'expérience et identifier les problèmes d'utilisabilité.
4. Analyser et agréger les données.
5. Présenter les résultats.

\textbf{Ressources :}

\begin{itemize}
\item Introduction à l'évaluation heuristique : un guide pour débutants \url{https://uxdesign.cc/introduction-to-heuristic-evaluation-658705606518}
\item Les 10 heuristiques d'utilisabilité de Nielsen \url{https://www.nngroup.com/articles/ten-usability-heuristics/}
\item Ce que vous retirez vraiment d'une évaluation heuristique \url{https://uxmag.com/articles/what-you-really-get-from-a-heuristic-evaluation}
\item Analyse heuristique pour l'UX : comment mener une évaluation d'utilisabilité \url{https://uxdesign.cc/heuristic-analysis-for-ux-how-to-run-a-usability-evaluation-12c86d43936f}
\item L'analyse heuristique dans le processus de design \url{https://uxdesign.cc/heuristic-analysis-in-the-design-process-usability-inspection-methods-d200768eb38d}
\item Une nouvelle checklist d'évaluation heuristique d'utilisabilité \url{https://uxplanet.org/a-new-usability-heuristic-evaluation-checklist-259f588da308}
\item Outil : checklist d'évaluation heuristique \url{https://drive.google.com/file/d/10KbfbNZA1oVS1sXbjjXLPPmdZ6nqVkdc/view}
\end{itemize}


\chapter{Exercice 01 : UI V - Ex 2 - Recommandations}
\exnumber{01}
\extitle{UI V - Ex 2 - Recommandations}
\exfiles{ex01/PDF, nomme Name\_FirstName\_DeliverableName\_Date\_VersionNumber}
\exauthorize{Ceux nommés dans l'exercice}
\makeheaderfiles

\textbf{Instructions :}

À partir de l'audit réalisé à l'Ex.1, rédigez un ensemble de 3 à 6 recommandations pour améliorer les heuristiques et l'utilisabilité.

\textbf{Sites ou applications :}

\begin{itemize}
\item https://www.ikea.com/ \url{https://www.ikea.com/}
\item https://www.doctorswithoutborders.org/ \url{https://www.doctorswithoutborders.org/}
\item Application WikiArt
\end{itemize}

\textbf{Livrable :}

\begin{itemize}
\item Un document PDF avec 3 à 6 recommandations
\item Déposez le PDF sur Github, nommé ainsi : « Name\_FirstName\_DeliverableName\_Date\_VersionNumber ». Exemple : « Doe\_John\_MidFiPrototype\_05242024\_V1 ».
\item Ajoutez un titre dans le document.
\item Soignez l'aspect global du document. Il doit être clair, simple et lisible. Vous pouvez regarder canva.com, sans surcharger le dessin de détails.
\end{itemize}

\textbf{Pistes :}

\begin{itemize}
\item Classez vos recommandations selon les critères heuristiques
\end{itemize}


\end{document}
